\documentclass[11pt]{article}

\usepackage[margin=1in]{geometry}
\usepackage{amsmath,amssymb,amsthm}
\usepackage{bm}
\usepackage{enumitem}
\usepackage{mathtools}
\usepackage{microtype}

% ---------- operators / macros ----------
\DeclareMathOperator{\tr}{tr}
\DeclareMathOperator{\rank}{rank}
\DeclareMathOperator{\Var}{Var}
\DeclareMathOperator{\Cov}{Cov}
\DeclareMathOperator{\Corr}{Corr}
\newcommand{\Cs}{\mathcal{C}}
\newcommand{\E}{\mathbb{E}}
\newcommand{\R}{\mathbb{R}}
\newcommand{\N}{\mathcal{N}}
\newcommand{\bX}{\mathbf{X}}
\newcommand{\bY}{\mathbf{Y}}
\newcommand{\bZ}{\mathbf{Z}}
\newcommand{\bH}{\mathbf{H}}
\newcommand{\bI}{\mathbf{I}}
\newcommand{\bP}{\mathbf{P}}
\newcommand{\bx}{\mathbf{x}}
\newcommand{\bu}{\mathbf{u}}
\newcommand{\bbeta}{\bm{\beta}}
\newcommand{\btheta}{\bm{\theta}}
\newcommand{\bmu}{\bm{\mu}}
\newcommand{\bSig}{\bm{\Sigma}}
\newcommand{\beps}{\bm{\varepsilon}}
\newcommand{\beeta}{\bm{\eta}}
\newcommand{\T}{^{\mathsf{T}}}
\newcommand{\inv}{^{-1}}
\newcommand{\Sxx}{S_{xx}}
\newcommand{\bhat}[1]{\widehat{#1}}

\theoremstyle{remark}
\newtheorem*{remark}{Remark}

\setlength{\parindent}{0pt}
\setlength{\parskip}{4pt}
\setlength{\emergencystretch}{3em}

\newcommand{\ans}{\medskip\noindent\textbf{Solution.}\ }

\begin{document}

\begin{center}
	{\large\bfseries Linear Statistical Models --- Homework 1}\\[2pt]
	{\normalsize B.Stat.\ 3rd year, Semester 1, 2026--27}\\[2pt]
	{\small Solutions}
\end{center}

\hrule
\medskip

\noindent\textbf{References to lecture notes.}\;
\textbf{[GM]} \emph{Gauss--Markov Framework for Linear Models};\;
\textbf{[SLR]} \emph{Overview of Simple Linear Regression};\;
\textbf{[PMI]} \emph{Projection and Matrix Inversion}.
Equation/result numbers below refer to these notes. Any result not stated there is proved.

\bigskip

% =====================================================================
\section*{Problem 1}

Let $X\sim\N_n(\bmu,\bSig)$ and let $A$ be a symmetric $n\times n$ matrix.

\medskip
\noindent\emph{A standardisation used in (b),(c).}
Since $\bSig$ is a covariance matrix it is symmetric positive semidefinite, so $\bSig=LL\T$ for some
$n\times n$ matrix $L$. If $U\sim\N_n(\mathbf 0,\bI_n)$, the affine image $\bmu+LU$ is Gaussian with mean
$\bmu$ and covariance $L\,\Var(U)\,L\T=LL\T=\bSig$; hence $\bmu+LU\sim\N_n(\bmu,\bSig)$. We may therefore
take
\[
X=\bmu+LU,\qquad U\sim\N_n(\mathbf 0,\bI_n).
\]
For i.i.d.\ standard normal coordinates $U_1,\dots,U_n$ we use the elementary moments
$\E U_i=0,\ \E U_i^2=1,\ \E U_i^3=0,\ \E U_i^4=3$ and independence, which give
\begin{align}
\E[U_iU_j]&=\delta_{ij}, \label{eq:m2}\\
\E[U_iU_jU_k]&=0, \label{eq:m3}\\
\E[U_iU_jU_kU_\ell]&=\delta_{ij}\delta_{k\ell}+\delta_{ik}\delta_{j\ell}+\delta_{i\ell}\delta_{jk}.
\label{eq:m4}
\end{align}
Indeed \eqref{eq:m3} holds because at least one index is unpaired (its factor contributes $\E U_i=0$), and
\eqref{eq:m4} is checked by cases: all four indices equal gives $\E U_i^4=3=1+1+1$; two distinct pairs gives
$1$; otherwise an unpaired index gives $0$ --- exactly what the right-hand side returns.

% ---------------------------------------------------------------
\subsection*{(a) $\E(X\T AX)=\tr(A\bSig)+\bmu\T A\bmu$}

This needs only $\E X=\bmu$ and $\Var(X)=\bSig$. From
$\bSig=\E[(X-\bmu)(X-\bmu)\T]=\E[XX\T]-\bmu\bmu\T$ we get $\E[XX\T]=\bSig+\bmu\bmu\T$. Since $X\T AX$ is a
scalar, $X\T AX=\tr(X\T AX)=\tr(AXX\T)$, and by linearity of $\tr$ and $\E$,
\[
\E(X\T AX)=\tr\!\big(A\,\E[XX\T]\big)=\tr\!\big(A(\bSig+\bmu\bmu\T)\big)
=\tr(A\bSig)+\tr(\bmu\T A\bmu)=\tr(A\bSig)+\bmu\T A\bmu. \qquad\blacksquare
\]

% ---------------------------------------------------------------
\subsection*{(b) $\Cov(X,X\T AX)=2\bSig A\bmu$}

Write $X=\bmu+LU$. Using symmetry of $A$ ($\bmu\T A L U=U\T L\T A\bmu$, both scalars),
\[
X\T AX=\bmu\T A\bmu+2\bmu\T ALU+U\T\!\underbrace{L\T AL}_{=:B}\,U .
\]
The constant $\bmu\T A\bmu$ does not affect covariance, and $X-\E X=LU$, so
\[
\Cov(X,X\T AX)=\Cov\!\big(LU,\;2(L\T A\bmu)\T U\big)+\Cov\!\big(LU,\;U\T BU\big).
\]
With $v:=L\T A\bmu$, the first term is $2L\,\E[UU\T]\,v=2Lv=2LL\T A\bmu=2\bSig A\bmu$ by \eqref{eq:m2}.
The $i$-th component of the second term is
$\sum_{j,k}B_{jk}\,\E[U_iU_jU_k]=0$ by \eqref{eq:m3}. Hence
$\Cov(X,X\T AX)=2\bSig A\bmu.$ \hfill$\blacksquare$

% ---------------------------------------------------------------
\subsection*{(c) $\Var(X\T AX)=2\tr\!\big((A\bSig)^2\big)+4\bmu\T A\bSig A\bmu$}

From (b), $X\T AX=\text{const}+2v\T U+U\T BU$ with $v=L\T A\bmu$, $B=L\T AL$. Thus
\[
\Var(X\T AX)=4\Var(v\T U)+\Var(U\T BU)+4\Cov(v\T U,\;U\T BU).
\]
\emph{Linear part.} $\Var(v\T U)=v\T\Var(U)v=v\T v=\bmu\T A LL\T A\bmu=\bmu\T A\bSig A\bmu$, so
$4\Var(v\T U)=4\bmu\T A\bSig A\bmu$.

\emph{Cross term.} $\Cov(v\T U,U\T BU)=\sum_i v_i\sum_{j,k}B_{jk}\,\E[U_iU_jU_k]=0$ by \eqref{eq:m3}.

\emph{Quadratic part.} $\E[U\T BU]=\sum_i B_{ii}=\tr B$, and by \eqref{eq:m4},
\[
\E\big[(U\T BU)^2\big]=\sum_{i,j,k,\ell}B_{ij}B_{k\ell}\,\E[U_iU_jU_kU_\ell]
=(\tr B)^2+\sum_{i,j}B_{ij}^2+\sum_{i,j}B_{ij}B_{ji}
=(\tr B)^2+2\tr(B^2),
\]
where the last equality uses $B\T=B$. Hence $\Var(U\T BU)=2\tr(B^2)$. Finally, since $LL\T=\bSig$,
\[
\tr(B^2)=\tr(L\T AL\,L\T AL)=\tr(A\,LL\T\!A\,LL\T)=\tr\!\big((A\bSig)^2\big).
\]
Collecting the three parts,
\[
\Var(X\T AX)=4\bmu\T A\bSig A\bmu+2\tr\!\big((A\bSig)^2\big). \qquad\blacksquare
\]

% =====================================================================
\section*{Problem 2}

$X\sim\N_n(\bmu,\bSig)$ with $\bSig$ positive definite; $A,B$ symmetric. \emph{Claim:} $X\T AX$ and $BX$
are independent $\iff B\bSig A=\mathbf 0$.

\ans
Because $\bSig$ is positive definite we may write $\bSig=LL\T$ with $L$ \emph{invertible} (e.g.\ a Cholesky
factor), and put $X=\bmu+LU$, $U\sim\N_n(\mathbf 0,\bI_n)$. Set
\[
M:=BL,\qquad P:=L\T AL\ (\text{symmetric}),\qquad a:=L\T A\bmu .
\]
Then $BX=B\bmu+MU$ and $X\T AX=c+2a\T U+U\T PU$ with $c=\bmu\T A\bmu$. Note
\[
MP=BL\,L\T AL=B\bSig A\,L,\qquad Ma=BL\,L\T A\bmu=B\bSig A\bmu,
\]
so, as $L$ is invertible, $MP=\mathbf 0\iff B\bSig A=\mathbf 0$, and this implies $Ma=B\bSig A\bmu=\mathbf 0$.

\medskip
\emph{Joint MGF.} For $z\sim\N_n(\mathbf 0,\bI_n)$ and any vector $b$ and symmetric $Q$ with $\bI-2Q$
positive definite,
\[
\E\exp\!\big(b\T z+z\T Qz\big)
=\int\frac{e^{-\frac12 z\T(\bI-2Q)z+b\T z}}{(2\pi)^{n/2}}\,dz
=\det(\bI-2Q)^{-1/2}\exp\!\Big(\tfrac12 b\T(\bI-2Q)\inv b\Big),
\]
by completing the square. Apply this with $s\in\R^n$, $t\in\R$ small (so $\bI-2tP\succ0$): since
$s\T BX+t\,X\T AX=(s\T B\bmu+tc)+(M\T s+2ta)\T U+t\,U\T PU$, writing $N:=(\bI-2tP)\inv$ and $b:=M\T s+2ta$,
\[
\varphi(s,t):=\E\,e^{\,s\T BX+t X\T AX}
=e^{\,s\T B\bmu+tc}\,\det(\bI-2tP)^{-1/2}\exp\!\big(\tfrac12 b\T N b\big).
\]
The marginals are $\varphi(s,0)=e^{s\T B\bmu}\exp(\tfrac12 s\T MM\T s)$ and
$\varphi(0,t)=e^{tc}\det(\bI-2tP)^{-1/2}\exp(2t^2 a\T N a)$. Since the MGF is finite on a neighbourhood of the
origin, $BX$ and $X\T AX$ are independent \emph{iff} $\varphi(s,t)=\varphi(s,0)\varphi(0,t)$ there, i.e.\ iff
$\tfrac12 b\T Nb=\tfrac12 s\T MM\T s+2t^2a\T Na$. Expanding $b\T Nb=s\T MNM\T s+4t\,s\T MNa+4t^2 a\T Na$ and
cancelling $2t^2a\T Na$, the condition becomes
\begin{equation*}
\tfrac12\,s\T MNM\T s+2t\,s\T MNa=\tfrac12\,s\T MM\T s\qquad\text{for all }s\in\R^n,\ t\ \text{near }0.
\tag{$\star$}
\end{equation*}
For fixed $t$, ($\star$) holds for all $s$ iff its quadratic and linear parts in $s$ match separately
(replace $s\mapsto-s$), i.e.
\[
M N M\T=MM\T\quad\text{and}\quad MNa=\mathbf 0,\qquad N=(\bI-2tP)\inv .
\]
Using the convergent expansion $N=\sum_{k\ge0}(2t)^kP^k$ for small $t$:
\[
MNM\T-MM\T=\sum_{k\ge1}(2t)^k MP^kM\T=\mathbf 0\ \forall\,t
\iff MP^kM\T=\mathbf0\ \forall k\ge1 .
\]
Taking $k=2$ gives $MP^2M\T=(MP)(MP)\T=\mathbf 0$, and $CC\T=\mathbf 0\Rightarrow C=\mathbf 0$ (its diagonal
entries are squared row norms); hence $MP=\mathbf 0$. Conversely $MP=\mathbf 0$ gives $MP^kM\T=\mathbf0$ for
all $k\ge1$ and $MN=M$, so also $MNa=Ma$. Thus ($\star$) holds for all $s,t$ iff
\[
MP=\mathbf 0\ \text{ and }\ Ma=\mathbf 0.
\]
Finally $MP=\mathbf0\iff B\bSig A=\mathbf0$, and $B\bSig A=\mathbf0$ already forces $Ma=B\bSig A\bmu=\mathbf0$.
Therefore
\[
X\T AX\ \text{and}\ BX\ \text{are independent}\iff B\bSig A=\mathbf 0. \qquad\blacksquare
\]

\begin{remark}
The ``if'' direction has a shorter argument: $B\bSig A=\mathbf0\Rightarrow A\bSig B=(B\bSig A)\T=\mathbf0$, so
$\Cov(AX,BX)=A\bSig B\T=\mathbf0$; jointly normal and uncorrelated $\Rightarrow AX\perp BX$; and
$X\T AX=(AX)\T A^{+}(AX)$ (using the Moore--Penrose inverse $A^{+}$, which exists by [PMI, Prop.~15], and
$AA^{+}A=A$) is a function of $AX$, hence independent of $BX$. The MGF computation is needed for ``only if''.
\end{remark}

% =====================================================================
\section*{Problem 3}

Model $Y_i=\beta_0+\beta_1x_i+\varepsilon_i$, $\bx_i=(1,x_i)\T$, $\bX$ the $n\times2$ design with rows
$\bx_i\T$, and $\bbeta=(\beta_0,\beta_1)\T$. Assume $\bX$ has full column rank, so $\bX\T\bX$ is nonsingular
and $\bhat{\bbeta}^{(n)}=(\bX\T\bX)\inv\bX\T\bY$ \;[GM, full-rank LSE]. Write $G:=(\bX\T\bX)\inv$
and $\bu:=\bx_{n+1}$.

\medskip
\noindent\emph{A rank-one inverse update (Sherman--Morrison).} For nonsingular $M$ and a vector $\bu$ with
$1+\bu\T M\inv\bu\ne0$,
\begin{equation*}
(M+\bu\bu\T)\inv=M\inv-\frac{M\inv\bu\bu\T M\inv}{1+\bu\T M\inv\bu}.
\tag{SM}
\end{equation*}
\emph{Proof:} multiply the right-hand side by $(M+\bu\bu\T)$ and let $\gamma:=1+\bu\T M\inv\bu$:
\begin{align*}
&(M+\bu\bu\T)\Big(M\inv-\tfrac{1}{\gamma}M\inv\bu\bu\T M\inv\Big)\\
&\qquad=\bI+\bu\bu\T M\inv-\tfrac1\gamma\bu\bu\T M\inv-\tfrac1\gamma\bu(\bu\T M\inv\bu)\bu\T M\inv\\
&\qquad=\bI+\bu\bu\T M\inv\Big(1-\tfrac{\gamma}{\gamma}\Big)=\bI. \quad\checkmark
\end{align*}
Here $M=\bX\T\bX$ is positive definite, so $\bu\T G\bu\ge0$ and $\gamma=1+\bu\T G\bu\ge1>0$; thus
$c_{n+1}=(1+\bx_{n+1}\T G\,\bx_{n+1})\inv\in(0,1]$ is well defined.

% ---------------------------------------------------------------
\subsection*{(a)}

Adding $(x_{n+1},Y_{n+1})$ stacks one row $\bx_{n+1}\T$ on $\bX$ and one entry $Y_{n+1}$ on $\bY$, so
\[
\bX_{n+1}\T\bX_{n+1}=\bX\T\bX+\bu\bu\T,\qquad
\bX_{n+1}\T\bY_{(n+1)}=\bX\T\bY+\bu\,Y_{n+1}.
\]
By (SM) with $M=\bX\T\bX$, and writing $c:=c_{n+1}$,
$(\bX_{n+1}\T\bX_{n+1})\inv=G-c\,G\bu\bu\T G$. Hence
\[
\bhat{\bbeta}^{(n+1)}=(G-c\,G\bu\bu\T G)(\bX\T\bY+\bu Y_{n+1})
=\bhat{\bbeta}^{(n)}+G\bu\,Y_{n+1}-c\,G\bu\,\bu\T\bhat{\bbeta}^{(n)}-c\,G\bu(\bu\T G\bu)Y_{n+1},
\]
using $G\bX\T\bY=\bhat{\bbeta}^{(n)}$ and $\bu\T G\bX\T\bY=\bu\T\bhat{\bbeta}^{(n)}$. Now
$\bu\T\bhat{\bbeta}^{(n)}=\bx_{n+1}\T\bhat{\bbeta}^{(n)}=\bhat Y^{(n)}_{n+1}$ (the value predicted for
$Y_{n+1}$ from the first $n$ observations --- this is the BLUE of
$\E Y_{n+1}=\bx_{n+1}\T\bbeta$, [GM, Thm~1]), and $c(\bu\T G\bu)=c(\tfrac1c-1)=1-c$. Factoring out $G\bu$,
\[
\bhat{\bbeta}^{(n+1)}-\bhat{\bbeta}^{(n)}
=G\bu\Big[\,Y_{n+1}\big(1-c\,\bu\T G\bu\big)-c\,\bhat Y^{(n)}_{n+1}\Big]
=G\bu\Big[c\,Y_{n+1}-c\,\bhat Y^{(n)}_{n+1}\Big],
\]
because $1-c\,\bu\T G\bu=1-(1-c)=c$. Therefore
\[
\boxed{\ \bhat{\bbeta}^{(n+1)}-\bhat{\bbeta}^{(n)}
=c_{n+1}\big(Y_{n+1}-\bhat Y^{(n)}_{n+1}\big)(\bX\T\bX)\inv\bx_{n+1}.\ } \qquad\blacksquare
\]

% ---------------------------------------------------------------
\subsection*{(b)}

Left-multiply the boxed identity by $\bx_{n+1}\T=\bu\T$. Since
$\bhat Y^{(n+1)}_{n+1}=\bx_{n+1}\T\bhat{\bbeta}^{(n+1)}$ and $\bhat Y^{(n)}_{n+1}=\bu\T\bhat{\bbeta}^{(n)}$,
\[
\bhat Y^{(n+1)}_{n+1}-\bhat Y^{(n)}_{n+1}
=c_{n+1}\big(Y_{n+1}-\bhat Y^{(n)}_{n+1}\big)\,\bu\T G\bu .
\]
Subtract both sides from $Y_{n+1}-\bhat Y^{(n)}_{n+1}$ and use $1-c_{n+1}\,\bu\T G\bu=c_{n+1}$:
\[
Y_{n+1}-\bhat Y^{(n+1)}_{n+1}
=\big(Y_{n+1}-\bhat Y^{(n)}_{n+1}\big)\big(1-c_{n+1}\bu\T G\bu\big)
=c_{n+1}\big(Y_{n+1}-\bhat Y^{(n)}_{n+1}\big).
\]
Dividing by $c_{n+1}\in(0,1]$,
\[
Y_{n+1}-\bhat Y^{(n)}_{n+1}=\frac{Y_{n+1}-\bhat Y^{(n+1)}_{n+1}}{c_{n+1}}. \qquad\blacksquare
\]

% =====================================================================
\section*{Problem 4}

Normal simple linear regression, $\varepsilon_i\overset{\text{iid}}{\sim}\N(0,\sigma^2)$; test
$H_0:\beta_1=0$ vs $H_1:\beta_1\ne0$. Put $\Sxx:=\sum_i(x_i-\bar x)^2$.

\begin{remark}
With two parameters $(\beta_0,\beta_1)$ the residual degrees of freedom are $n-2$. The correct statement ---
matching [SLR, eq.~(12)] and part (b) below --- is $T\sim t_{n-2}$; the ``$t_{(n-1)}$'' in the problem is a
typo for $t_{(n-2)}$. We prove $T\sim t_{n-2}$ under $H_0$.
\end{remark}

% ---------------------------------------------------------------
\subsection*{(a) $T=\bhat\beta_1/s(\bhat\beta_1)\sim t_{n-2}$ under $H_0$}

Write $\bhat\beta_1=\sum_i c_iY_i$ with $c_i=(x_i-\bar x)/\Sxx$ \;[SLR, eq.~(4)]. Then
$\sum_i c_i=0$ and $\sum_i c_ix_i=\sum_i(x_i-\bar x)x_i/\Sxx=1$, so
$\E\bhat\beta_1=\sum_i c_i(\beta_0+\beta_1x_i)=\beta_1$ and
$\Var(\bhat\beta_1)=\sigma^2\sum_i c_i^2=\sigma^2/\Sxx$ \;[SLR, eq.~(7)]. Being a linear combination of
independent normals, $\bhat\beta_1$ is normal:
\[
\bhat\beta_1\sim\N\!\Big(\beta_1,\ \tfrac{\sigma^2}{\Sxx}\Big),
\qquad\text{so under }H_0:\ \ Z:=\frac{\bhat\beta_1\sqrt{\Sxx}}{\sigma}\sim\N(0,1).
\]
By [SLR, eq.~(12)], $\mathrm{SSE}\sim\sigma^2\chi^2_{n-2}$ and $\mathrm{SSE}$ is independent of $\bhat\beta_1$;
write $W:=\mathrm{SSE}/\sigma^2\sim\chi^2_{n-2}$, independent of $Z$. Since $s^2(\bhat\beta_1)=\mathrm{MSE}/\Sxx$
with $\mathrm{MSE}=\mathrm{SSE}/(n-2)$ \;[SLR, eq.~(8)],
\[
T=\frac{\bhat\beta_1}{s(\bhat\beta_1)}
=\frac{\bhat\beta_1\sqrt{\Sxx}}{\sqrt{\mathrm{MSE}}}
=\frac{\bhat\beta_1\sqrt{\Sxx}/\sigma}{\sqrt{\mathrm{SSE}/\big((n-2)\sigma^2\big)}}
=\frac{Z}{\sqrt{W/(n-2)}}.
\]
A ratio $Z/\sqrt{W/k}$ with $Z\sim\N(0,1)\perp W\sim\chi^2_k$ is, by definition, $t_k$. Hence under $H_0$,
$T\sim t_{n-2}$. \hfill$\blacksquare$

% ---------------------------------------------------------------
\subsection*{(b) $F=\mathrm{MSR}/\mathrm{MSE}=T^2$}

The ANOVA decomposition $\mathrm{SSTO}=\mathrm{SSR}+\mathrm{SSE}$ has $\mathrm{d.f.}(\mathrm{SSR})=1$ with
$\mathrm{SSR}=\bhat\beta_1^{\,2}\,\Sxx$ \;[SLR, ANOVA]. Thus
$\mathrm{MSR}=\mathrm{SSR}/1=\bhat\beta_1^{\,2}\,\Sxx$ and, since $s^2(\bhat\beta_1)=\mathrm{MSE}/\Sxx$,
\[
F=\frac{\mathrm{MSR}}{\mathrm{MSE}}=\frac{\bhat\beta_1^{\,2}\,\Sxx}{\mathrm{MSE}}
=\frac{\bhat\beta_1^{\,2}}{\mathrm{MSE}/\Sxx}
=\frac{\bhat\beta_1^{\,2}}{s^2(\bhat\beta_1)}
=\Big(\frac{\bhat\beta_1}{s(\bhat\beta_1)}\Big)^{2}=T^2. \qquad\blacksquare
\]

% ---------------------------------------------------------------
\subsection*{(c) The two tests always agree}

If $T\sim t_{n-2}$ then $T^2=Z^2/(W/(n-2))$ with $Z^2\sim\chi^2_1\perp W\sim\chi^2_{n-2}$, so
$T^2\sim F_{1,n-2}$ (consistent with (b)). For the level-$\alpha$ tests:
\[
\text{$t$-test rejects}\iff |T|>t_{n-2}(1-\tfrac\alpha2),\qquad
\text{$F$-test rejects}\iff F=T^2>F_{1,n-2}(1-\alpha).
\]
Let $c:=t_{n-2}(1-\tfrac\alpha2)$. Then
\[
\Pr\big(F_{1,n-2}>c^2\big)=\Pr(T^2>c^2)=\Pr(|T|>c)
=\Pr(T>c)+\Pr(T<-c)=\tfrac\alpha2+\tfrac\alpha2=\alpha,
\]
so $c^2$ is exactly the upper-$\alpha$ point of $F_{1,n-2}$, i.e.\ $F_{1,n-2}(1-\alpha)=c^2$. Consequently
\[
|T|>c\iff T^2>c^2\iff F>F_{1,n-2}(1-\alpha),
\]
so for every data set and every $\alpha$ the two rejection regions coincide; the tests reach identical
conclusions. \hfill$\blacksquare$

% =====================================================================
\section*{Problem 5}

$\bH=\bX(\bX\T\bX)\inv\bX\T$ with rows $\bx_i\T=(1,x_i)$; $h_{ii}=\bx_i\T(\bX\T\bX)\inv\bx_i$ is the $i$-th
diagonal of $\bH$. Let $\bhat Y_i=\bx_i\T\bhat{\bbeta}$ (full-data fit) and $\bhat Y_i^{(-i)}$ the fit at
$x_i$ from the regression that omits observation $i$. \emph{Claim:}
$Y_i-\bhat Y_i^{(-i)}=(Y_i-\bhat Y_i)/(1-h_{ii})$.

\ans
Apply Problem 3 with the roles reversed: take the ``base'' data set to be the $n-1$ observations
$\{(x_j,Y_j)\}_{j\ne i}$ (design $\bX_{(-i)}$, assumed full rank) and let observation $i$ play the role of the
``$(n{+}1)$-th'' added point. In Problem 3's notation this identifies
\[
\bhat{\bbeta}^{(n)}\!\to\bhat{\bbeta}^{(-i)},\quad
\bhat{\bbeta}^{(n+1)}\!\to\bhat{\bbeta},\quad
\bhat Y^{(n)}_{n+1}\!\to\bhat Y_i^{(-i)},\quad
\bhat Y^{(n+1)}_{n+1}\!\to\bhat Y_i,\quad
c_{n+1}\!\to c:=\big(1+\bx_i\T M\inv\bx_i\big)\inv,
\]
where $M:=\bX_{(-i)}\T\bX_{(-i)}$. Problem 3(b) then reads
\begin{equation*}
Y_i-\bhat Y_i^{(-i)}=\frac{Y_i-\bhat Y_i}{c}.
\tag{5.1}
\end{equation*}
It remains to show $c=1-h_{ii}$. Deleting row $i$ gives $\bX\T\bX=M+\bx_i\bx_i\T$, so with
$g:=\bx_i\T M\inv\bx_i$ and (SM) from Problem 3,
\[
h_{ii}=\bx_i\T(\bX\T\bX)\inv\bx_i
=\bx_i\T M\inv\bx_i-\frac{(\bx_i\T M\inv\bx_i)^2}{1+\bx_i\T M\inv\bx_i}
=g-\frac{g^2}{1+g}=\frac{g}{1+g}.
\]
Therefore $1-h_{ii}=\dfrac{1}{1+g}=\big(1+\bx_i\T M\inv\bx_i\big)\inv=c$. Substituting into (5.1),
\[
Y_i-\bhat Y_i^{(-i)}=\frac{Y_i-\bhat Y_i}{1-h_{ii}},\qquad i=1,\dots,n. \qquad\blacksquare
\]

% =====================================================================
\section*{Problem 6}

$\mathcal M_1:\bY=\bX\bbeta+\beps$ and $\mathcal M_2:\bY=\bZ\btheta+\beeta$, both with mean-zero errors of
covariance $\sigma^2\bI$, $\bZ$ of full column rank. Being a \emph{reparametrisation} means the two models
describe the same mean vector, i.e.\ $\{\bX\bbeta\}=\{\bZ\btheta\}$, so the model spaces coincide:
$\Cs(\bX)=\Cs(\bZ)=:\mathcal M$.

\ans
\emph{BLUE of $\bX\bbeta$.} Each component $(\bX\bbeta)_i=\bx_i\T\bbeta$ has coefficient vector $\bx_i$ (the
$i$-th row of $\bX$) lying in $\Cs(\bX\T)=\mathcal R(\bX)$, hence is estimable [GM, Lemma~1]; by the
Gauss--Markov theorem [GM, Thm~1] its BLUE is the corresponding component of the (unique) fitted vector
\[
\widehat{\bX\bbeta}=\bP_{\!\mathcal M}\,\bY,
\]
the orthogonal projection of $\bY$ onto $\mathcal M=\Cs(\bX)$ [GM, eq.~(2)]. The orthogonal projector onto a
subspace is unique [PMI, Prop.~4], so it is the same whether computed from $\bX$ or $\bZ$:
$\bP_{\!\mathcal M}=\bX(\bX\T\bX)\inv\bX\T=\bZ(\bZ\T\bZ)\inv\bZ\T$ (the latter valid since $\bZ$ has full
column rank, [PMI, Prop.~5]).

\emph{BLUE of $\btheta$.} As $\bZ$ has full column rank, every $\btheta_j=\mathbf e_j\T\btheta$ is estimable
and, by the Gauss--Markov theorem, the BLUE of $\btheta$ is the least-squares estimator under $\mathcal M_2$,
\[
\bhat\btheta=(\bZ\T\bZ)\inv\bZ\T\bY .
\]
Moreover $\bZ\bhat\btheta=\bZ(\bZ\T\bZ)\inv\bZ\T\bY=\bP_{\!\mathcal M}\bY$.

\emph{The two BLUEs in terms of each other.} Combining the displays,
\[
\boxed{\ \widehat{\bX\bbeta}=\bP_{\!\mathcal M}\bY=\bZ\bhat\btheta\ }
\qquad\text{and}\qquad
\boxed{\ \bhat\btheta=(\bZ\T\bZ)\inv\bZ\T\,\widehat{\bX\bbeta}.\ }
\]
The second identity is consistent since, using $\bP_{\!\mathcal M}\bY=\bZ\bhat\btheta$,
\begin{align*}
(\bZ\T\bZ)\inv\bZ\T\bP_{\!\mathcal M}\bY
&=(\bZ\T\bZ)\inv\bZ\T\bZ(\bZ\T\bZ)\inv\bZ\T\bY\\
&=(\bZ\T\bZ)\inv\bZ\T\bY=\bhat\btheta.
\end{align*}
Thus the BLUE of $\bX\bbeta$ is $\bZ$ times the
BLUE of $\btheta$, and the BLUE of $\btheta$ is recovered from that of $\bX\bbeta$ by the left inverse
$(\bZ\T\bZ)\inv\bZ\T$ of $\bZ$. \hfill$\blacksquare$

\end{document}
